\begin{lemma}[Continuity of the horizontal gradient along a convergent pair]
Let $n,m\in\mathbb N$ and let
$\sigma:\mathbb R^{n}\to(\mathbb R^{n})^{m}$, written $x\mapsto(\sigma^{1}(x),\ldots,
\sigma^{m}(x))$, be a frame each of whose components $x\mapsto\sigma^{i}(x)$ is continuous.
Let $x_k\to x_\infty$ in $\mathbb R^{n}$, and let $D_k\to D_\infty$ in the operator norm,
where each $D_k$ and $D_\infty$ is a continuous linear functional on $\mathbb R^{n}$. Then
\[
 \mathrm{hgrad}\ \sigma\ x_k\ D_k\ \longrightarrow\
 \mathrm{hgrad}\ \sigma\ x_\infty\ D_\infty \qquad\text{in }\mathbb R^{m},
\]
with $\mathrm{hgrad}$ as in \noderef{hgrad}.
\end{lemma}
\begin{proof}
(The frame is taken bare here, with continuity of its components as the only hypothesis;
in the application the $\sigma^{i}$ are the horizontal fields of a Carnot setting and are
continuous because they are smooth, but no group structure enters this argument.)

\medskip\noindent\textbf{Step 0: the components of $\mathrm{hgrad}$.} By \noderef{hgrad},
for every $y\in\mathbb R^{n}$ and every continuous linear functional $D$ on $\mathbb R^{n}$,
the vector $\mathrm{hgrad}\ \sigma\ y\ D\in\mathbb R^{m}$ is the one whose $i$-th component
is $D(\sigma^{i}(y))$. So the assertion is that, in $\mathbb R^{m}$,
\[
 \big(D_k(\sigma^{i}(x_k))\big)_{i}\ \longrightarrow\
 \big(D_\infty(\sigma^{i}(x_\infty))\big)_{i} . \tag{1}
\]

\medskip\noindent\textbf{Step 1: the jets converge uniformly.} We have
\[
 \|D_k-D_\infty\|\ \longrightarrow\ 0 . \tag{2}
\]
Indeed $D_k\to D_\infty$, and the constant sequence $D_\infty$ converges to $D_\infty$, so
$D_k-D_\infty\to D_\infty-D_\infty=0$; the norm is continuous, so
$\|D_k-D_\infty\|\to\|0\|=0$.

\medskip\noindent\textbf{Step 2: convergence of each component.} Fix $i$. We claim
\[
 D_k(\sigma^{i}(x_k))\ \longrightarrow\ D_\infty(\sigma^{i}(x_\infty)). \tag{3}
\]
First, $\sigma^{i}(x_k)\to\sigma^{i}(x_\infty)$: the map $x\mapsto\sigma^{i}(x)$ is
continuous by hypothesis, hence continuous at $x_\infty$, and $x_k\to x_\infty$, so the
composition converges. Consequently:
\begin{itemize}
\item $D_\infty(\sigma^{i}(x_k))\to D_\infty(\sigma^{i}(x_\infty))$, because $D_\infty$ is
  a continuous linear functional and is therefore continuous at $\sigma^{i}(x_\infty)$.
\item $\|D_k-D_\infty\|\,\|\sigma^{i}(x_k)\|\to0$: by $(2)$ the first factor tends to $0$,
  and the second converges to $\|\sigma^{i}(x_\infty)\|$ by continuity of the norm, so the
  product tends to $0\cdot\|\sigma^{i}(x_\infty)\|=0$.
\end{itemize}
Now compare the two sequences termwise. For each $k$, linearity gives
$D_\infty(\sigma^{i}(x_k))-D_k(\sigma^{i}(x_k))=(D_\infty-D_k)(\sigma^{i}(x_k))$, so by the
defining bound of the operator norm and $\|{-w}\|=\|w\|$,
\[
 \Big|D_\infty(\sigma^{i}(x_k))-D_k(\sigma^{i}(x_k))\Big|
  =\Big|(D_\infty-D_k)(\sigma^{i}(x_k))\Big|
  \le\|D_\infty-D_k\|\,\|\sigma^{i}(x_k)\|
  =\|D_k-D_\infty\|\,\|\sigma^{i}(x_k)\| .
\]
The left-hand side is the distance between the $k$-th terms of the two sequences; it is
non-negative and, by the second bullet, dominated by a sequence tending to $0$. Squeezing,
that distance tends to $0$. A sequence at vanishing distance from a convergent sequence has
the same limit, so from the first bullet we get $(3)$.

\medskip\noindent\textbf{Step 3: from components to $\mathbb R^{m}$.} Convergence in the
product topology on functions $\{1,\ldots,m\}\to\mathbb R$ is exactly convergence of each
component, so by $(3)$, applied to every $i$,
\[
 \big(i\mapsto D_k(\sigma^{i}(x_k))\big)\ \longrightarrow\
 \big(i\mapsto D_\infty(\sigma^{i}(x_\infty))\big)
\]
in that topology. The identification of such functions with $\mathbb R^{m}$ is a linear
homeomorphism, hence continuous, and composing a convergent sequence with a continuous map
preserves the limit. This gives $(1)$, which by Step 0 is the assertion of the lemma.
\end{proof}
